\documentclass{book}
\usepackage[utf8]{inputenc}
\usepackage[T1]{fontenc}
\usepackage{geometry}
\usepackage{graphicx}
\usepackage{float}
\usepackage{hyperref}
\usepackage[english,indonesian]{babel}

\title{Optimasi Pemrosesan Template LaTeX Menggunakan Pendekatan Berbasis AI}
\author{Academic Support Agent}
\date{\today}

\begin{document}
\maketitle

\begin{abstract}

\end{abstract}

\chapter{PENDAHULUAN}


\chapter{TINJAUAN PUSTAKA}


% Original content (commented out):
% \documentclass[
twoside,                % A una columna o a dos
openright,              % Los capítulos y secciones empiezan siempre a la derecha
12pt,                   % Tamaño de fuente
]{book}

\usepackage[spanish]{babel}         % Para el idioma español
\usepackage{eso-pic}                 % Para añadir elementos gráficos en cualquier lugar del documento
\usepackage[T1]{fontenc}             % Para la codificación de fuentes
\usepackage{lmodern} \normalfont     % Usa una fuente moderna (que no esté pixelada)
\DeclareFontShape{T1}{lmr}{bx}{sc} { <-> ssub * cmr/bx/sc }{}
\usepackage[doublespacing]{setspace} % Para establecer el espaciado doble
\usepackage{indentfirst}             % Para indentar el primer párrafo de cada sección
\usepackage{geometry}                % Para establecer los márgenes
\usepackage{graphicx}                % Para incluir gráficos
\usepackage{titlesec}                % Para personalizar los títulos de las secciones
\usepackage{fancyhdr}                % Para personalizar las cabeceras y pies de página
\usepackage{lipsum}                  % Para texto de relleno
\usepackage{amsmath,bm}              % Para símbolos y entornos matemáticos y negrita en matemáticas
\usepackage{minted}                  % Para resaltar código fuente
\usepackage{subcaption}              % Para subfiguras
\usepackage{mfirstuc}                % Para capitalizar la primera letra de palabras
\usepackage{appendix}                % Para gestionar apéndices
\usepackage{titletoc}                % Para personalizar la tabla de contenidos
\usepackage{emptypage}               % Para evitar números de página en páginas vacías
\usepackage{lscape}                  % Para poder poner páginas en horizontal
\usepackage{afterpage}               % Para ejecutar comandos después de la página actual
\usepackage{ragged2e}                % Para justificar el texto
\usepackage[style=apa,backend=biber]{biblatex} % Para la bibliografía en estilo APA
\addbibresource{referencias.bib}     % Añadir el archivo de bibliografía
\input{configuracion/datos_base.tex} % Incluir archivo de configuración
\graphicspath{{./figuras/}}          % Ruta para las figuras
\usepackage{csquotes}                % Para citas
\usepackage{hyperref}                % Para enlaces y referencias dentro del documento
\usepackage{url}                     % Para manejar URLs largas
\usepackage{breakurl}                % Opcional: Para dividir URLs en múltiples líneas
\usepackage{pdfpages}                % Para incluir PDF en las páginas
\usepackage{array}                   % Para mejorar tablas
\usepackage{booktabs}                % Para poder hacer tablas 
\usepackage{adjustbox}               % Paquete para ajustar contenido

\pagenumbering{gobble} % supress page numbers
\renewcommand{\rmdefault}{phv}
\renewcommand{\sfdefault}{phv} 

\usepackage{tikz}
\usepackage{afterpage}
\usepackage{pdflscape}

\usepackage[acronym,automake]{glossaries-extra} % Para el glosario y acrónimos
\setabbreviationstyle[acronym]{long-short}
\makeglossaries
\input{glossary}                     % Incluir el archivo de glosario
\let\cite\parencite
\let\citet\textcite
\geometry{
    top=2.54cm,
    left=2.54cm,
    right=2.54cm,
    bottom=2.54cm
}

% Configuración básica de la cabecera y pie de página
\pagestyle{fancy}
\fancyhf{} % Limpia cualquier configuración previa
\fancyhead[R]{\thepage} % Coloca el número de página en la parte superior derecha
\fancyhead[L]{\shorttitle} % Coloca el título corto en la parte superior izquierda

% Redefinir el estilo 'plain' para incluir el número de página en la parte superior derecha
\fancypagestyle{plain}{
  \fancyhf{}
  \fancyhead[R]{\thepage}
  \fancyhead[L]{\shorttitle}
  \fancyfoot{}
}

% Estilo cabecera con logos
\fancypagestyle{specialheader}{
  \fancyhf{}
  \fancyhead[L]{\includegraphics[height=2cm]{figuras/facfilylet.png}}
  \fancyhead[R]{\includegraphics[height=2cm]{figuras/EscudoUA.png}}
  \fancyfoot{} % No número de página
}

% Configuración de títulos de capítulo para "X. Título del capítulo" alineados a la izquierda y en negrita
\titleformat{\chapter}[hang]
{\LARGE\bfseries}  % Formato del título del capítulo: tamaño de fuente grande y negrita
{\thechapter.}  % Número del capítulo seguido de un punto
{1em}  % Espacio entre el número del capítulo y el título
{\LARGE}  % Tamaño del título del capítulo
[{\vspace{3ex}}]

% Configuración de títulos del preámbulo (secciones no numeradas)
\titleformat{\section}[hang]
{\LARGE\bfseries}  % Formato del título de la sección: tamaño de fuente grande y negrita
{}  % Sin número de sección
{0pt}  % Sin espacio adicional
{\LARGE}  % Tamaño del título de la sección
[{\vspace{3ex}}]

\begin{document}
% Portada y cubierta sin numeración
\pagestyle{empty}
\include{preliminares/cubierta}
\include{preliminares/portada}

% Páginas preliminares
\frontmatter
\pagestyle{plain}
\include{preliminares/dedicatoria}
\include{preliminares/agradecimientos}
\renewcommand{\contentsname}{Índice de Contenido}
\doublespacing
\titlecontents{chapter}[0pt]{\addvspace{1pc}\bfseries}{}{}{\hfill\contentspage}

% Indice del trabajo
\tableofcontents  % página para el índice del trabajo
\singlespacing
\newpage

% Indice de figuras
\renewcommand{\listfigurename}{Índice de Figuras}
\listoffigures   % página con el índice de figuras
\singlespacing  
\newpage

% Indice de tablas
\renewcommand{\listtablename}{Índice de tablas}
\renewcommand{\tablename}{Tabla}
\listoftables  % página con el índice de tablas
\singlespacing
\newpage

% Páginas con resúmenes
\include{preliminares/resumenES}
\include{preliminares/resumenEN}

% Aquí debería ir el glosario
\setglossarystyle{list}  % Writes the defined term in boldface font
\printglossary[type=main, title={Glosario}]
\newpage

% Aquí debería ir la lista de acrónimos
\setglossarystyle{list}  % Writes the defined term in boldface font
\printglossary[type=acronym, title={Lista de Acrónimos}]
\newpage


% Páginas de contenido
\mainmatter
\pagestyle{fancy}
\setcounter{chapter}{0}

% Ejemplo de uso de un acrónimo en el texto
\include{cuerpo/introduccion}
\include{cuerpo/estudioviabilidad}
\include{cuerpo/objetivos}
\include{cuerpo/estadodelarte}
\include{cuerpo/planificacion}
\include{cuerpo/metodologia}
\include{cuerpo/desarrollo}
\include{cuerpo/validacion}
\include{cuerpo/resultadosyconclusiones}

% Añadir la bibliografía al índice de contenido
\cleardoublepage
\addcontentsline{toc}{chapter}{Referencias}
\printbibliography[title={Referencias}]

% Páginas finales
\backmatter
\appendix
\clearpage
\include{apendices/apendiceA}
\include{apendices/leanstartupcanvas}

\end{document}


\end{document}